% ---------------------------------------------------------------------------
% Kapitel 3: Grundlagen und verwandte Arbeiten
% Wortgleich uebernommen aus der Confluence-Seite
% "Kapitel 3: Grundlagen und verwandte Arbeiten (Entwurf)"
% (Space BEH, Seiten-ID 34373645), Stand 04.10.2026.
% Abschnitt 3.5 Zusammenfassung steht auf der Seite noch nicht.
% ---------------------------------------------------------------------------

\chapter{Grundlagen und verwandte Arbeiten}
\label{kap:grundlagen}

Drei Stränge der Forschungsliteratur tragen diese Arbeit, und sie beantworten verschiedene Fragen. Die E-Portfolio-Literatur sagt, was ein Portfolio ist und woraus es besteht; die Akzeptanzforschung, warum ein freiwillig gefülltes System an seiner Bedienbarkeit scheitern kann; die Usability-Literatur, woran sich Bedienbarkeit bemisst. Keiner der Stränge leistet, was die beiden anderen leisten, und diese Arbeitsteilung ist selbst ein Ergebnis: Sie bestimmt, welche Anforderung sich auf einen Beleg stützen kann und welche eine eigene Ableitung bleibt. Das Kapitel bearbeitet damit die erste Teilfrage und liefert in Abschnitt 3.3 die Kriterien der zweiten.

\section{E-Portfolio: Begriffe und Kategorien}
\label{sec:eportfolio-begriffe}

Lorenzo und Ittelson (2005) bestimmen ein E-Portfolio als digitalisierte Sammlung von Artefakten, die eine Einzelperson, eine Gruppe, eine Gemeinschaft, eine Organisation oder eine Institution repräsentiert und aus Text und Grafik ebenso wie aus multimedialen Bestandteilen bestehen kann. Entscheidend für das vorliegende Vorhaben ist der Absatz, mit dem sie fortfahren: Ein E-Portfolio sei mehr als eine bloße Sammlung, es könne zugleich als Verwaltungswerkzeug dienen, das die enthaltenen Arbeiten ordnet und steuert, wer sie sehen und besprechen darf. Sichtbarkeitssteuerung ist damit im Begriff angelegt und keine Zutat, die erst beim Bau hinzukommt.

Aus der Beschreibung der Verwendungen gruppieren Lorenzo und Ittelson drei breite Kategorien: studentische, lehrbezogene und institutionelle Portfolios. Sie halten fest, dass zahlreiche Modelle diese Kategorien und die von ihnen genannten sechs Hauptfunktionen verbinden; die Einteilung ist also keine Typologie mit scharfen Rändern. Für die Einordnung des Vorhabens zählt eine Aussage, die beiläufig im Kategorienabschnitt fällt: Auf der breitesten Ebene lassen sich studentische und lehrbezogene Portfolios zu einem institutionellen zusammenführen, das Belege für Selbstbericht und Akkreditierung liefert. Lorenzo und Ittelson führen die drei Kategorien also nebeneinander und beschreiben die dritte zugleich als Aggregat der beiden anderen. Die vorliegende Arbeit liest sie vorrangig als Aggregationsebene; das ist eine Lesart und keine Festlegung der Quelle. Ein Portal, das studentische Arbeiten sammelt und nach außen zeigt, ist danach beides: der Ort, an dem sie entstehen, und das Aggregat, das die Einrichtung vertritt.

Für das institutionelle Portfolio geben Lorenzo und Ittelson eine Bestimmung von Kahn wieder, nach der es eine fokussierte Auswahl authentischer Arbeiten, Daten und Analysen ist, die Rechenschaft der Einrichtung belegt und zugleich der einrichtungsweiten Reflexion dient (Kahn, 2001, zitiert nach Lorenzo \& Ittelson, 2005). Sie nennt damit einen Zweck nach außen und einen nach innen. Zwei Folgerungen ergeben sich daraus. Erstens verlangt eine fokussierte Auswahl einen Auswahlschritt; eine Vollsammlung wäre etwas anderes. Zweitens zählt für ein Portal, das die Hochschule vertreten soll, die äußere Hälfte. Für die veröffentlichten Unterlagen einer Einrichtung beschreiben Lorenzo und Ittelson die Leserschaft als offen, bis hin zu jeder Person mit Internetzugang. Dass dasselbe für die darin gesammelten studentischen Arbeiten gelten soll, ist eine Übertragung dieser Arbeit; aus ihr folgt die Festlegung des Zugangs ohne Anmeldung. Ebenso eine Setzung ist die Einordnung als Präsentationsportfolio: Die Gegenüberstellung von Präsentations- und Lernportfolio findet sich bei Lorenzo und Ittelson nicht.

Was ein hochschulisches E-Portfolio enthalten sollte, beantwortet eine zweite Quelle. Mahasneh (2020) entwickelt aus einer Durchsicht erziehungswissenschaftlicher Studien und aus Interviews mit 20 Fachleuten der Erziehungswissenschaft ein Modell aus zehn Komponenten: Biografie der Studierenden, Kursplan, Berichte und Forschungsarbeiten, Hausaufgaben, Projekte und Experimente, Aktivitäten, Zusammenfassungen und Schlussfolgerungen, wissenschaftliches Material, Audio- und Videoclips sowie Proben studentischer Leistung. Geprüft wurde das Modell anschließend mit 90 Studierenden über einen Fragebogen aus 17 Aussagen auf fünfstufiger Skala; der Gesamtmittelwert der Einstellung beträgt 4,58.

Die zehn Komponenten lassen sich nach ihrem Bezugspunkt ordnen; diese Sortierung ist eine Leistung dieser Arbeit, denn Mahasneh führt sie unsortiert nebeneinander. Personenbezogen ist allein die Biografie. An eine laufende Lehrveranstaltung gebunden sind Kursplan, Hausaufgaben und Aktivitäten. Artefaktbezogen und von einer Veranstaltung unabhängig sind die übrigen sechs. Daraus folgt die Zweiteilung, die das Vorhaben trägt: eine personenbezogene Ebene über die Studienzeit hinweg und eine leistungsbezogene Ebene mit den einzelnen Arbeiten. Die drei veranstaltungsgebundenen Komponenten entfallen mit der in Abschnitt 1.2 gezogenen Abgrenzung gegenüber dem Lernmanagementsystem.

Zwei Einzelbefunde sind aussagekräftiger als der Gesamtwert. Alle 17 Aussagen ergänzen denselben Satzanfang, nach dem ein E-Portfolio bei etwas hilft. Die Aussage, es helfe beim Veröffentlichen eigener Erfahrungen zum Nutzen anderer, erreicht 4,96 und damit den zweiten Rang, bei der zweitgeringsten Streuung des Fragebogens; öffentliche Sichtbarkeit ist unter den Befragten also nahezu unstrittig. Über die Bewertung durch eine lehrende Person gehen die Antworten dagegen weiter auseinander als über jede andere Aussage außer einer, bleiben mit 4,38 aber im positiven Bereich. Fremdbewertung ist nicht abgelehnt, aber deutlich weniger einhellig getragen als die Veröffentlichung eigener Erfahrungen.

Die Reichweite beider Quellen ist begrenzt, und das gehört an dieselbe Stelle wie ihr Ertrag. Lorenzo und Ittelson (2005) legen einen Übersichtsbericht vor, der sich neben schriftlichen Quellen zu einem großen Teil auf Gespräche mit 34 Fachleuten stützt; seine Zahlen sind Momentaufnahmen von 2005. Mahasneh (2020) erhebt ohne Kontrollgruppe und ohne Vormessung an einem einzigen College in Jordanien. Die positiven Werte führt sie selbst auf die Neuartigkeit der Untersuchung, die erste Erfahrung der Befragten mit einem E-Portfolio und die im Unterricht sonst fehlende Gelegenheit zurück, Leistungen, Kreativität und Selbstbewertung zu zeigen. Das Modell trägt deshalb als Inhaltsraster und nicht als Wirkungsnachweis.

Lorenzo und Ittelson benennen daneben ein Problem, das in die folgenden Abschnitte führt. Sie zählen die Überflutung Lesender mit ungeordneter Information zu den Herausforderungen lehrbezogener Portfolios; dass sie ein Portal mit studentischen Arbeiten ebenso trifft, ist eine Übertragung dieser Arbeit. Der Fragenkatalog zur Einführung solcher Systeme, den sie aus einem Whitepaper des ePortConsortium übernehmen, fragt zudem, ob ein System benutzerfreundlich genug für eine leichte Übernahme sein werde (ePortConsortium, 2003, zitiert nach Lorenzo \& Ittelson, 2005). Die E-Portfolio-Literatur stellt die Akzeptanzfrage also selbst, beantwortet sie aber nicht.

\section{Technologieakzeptanzmodell (TAM)}
\label{sec:tam}

Davis (1989) bestimmt zwei Größen, mit denen sich die Annahme eines Systems erklären lässt. Die wahrgenommene Nützlichkeit ist der Grad, zu dem eine Person glaubt, dass die Benutzung eines Systems ihre Arbeitsleistung steigern würde. Die wahrgenommene Benutzerfreundlichkeit ist der Grad, zu dem sie diese Benutzung für frei von Anstrengung hält; Anstrengung gilt ihm dabei als endliche Ressource, die eine Person auf ihre Tätigkeiten verteilt. Unter sonst gleichen Bedingungen wird demnach die als leichter bedienbar wahrgenommene Anwendung eher angenommen.

Die empirische Grundlage sind zwei Untersuchungen mit zusammen 152 Personen und vier Anwendungen: eine Felduntersuchung mit 112 auswertbaren Antworten bei 120 ausgegebenen Fragebogen und eine Laboruntersuchung mit 40 Teilnehmenden nach einer Stunde Erprobung. Die beiden Skalen aus je sechs Aussagen erreichen Reliabilitäten von 0,98 und 0,94. In beiden Untersuchungen hängt die Nützlichkeit stärker mit der Nutzung zusammen als die Benutzerfreundlichkeit, mit 0,63 gegenüber 0,45 für die selbst berichtete gegenwärtige und mit 0,85 gegenüber 0,59 für die selbst vorhergesagte künftige Nutzung. Normwerte berichtet Davis nicht; der Punkt, an dem ein gemessener Wert eingeordnet wird, ist deshalb eine Festlegung dieser Arbeit und in Abschnitt 1.2 vorab benannt.

Der entscheidende Befund liegt nicht in der Stärke dieser Zusammenhänge, sondern in ihrer Anordnung. Sobald die Nützlichkeit statistisch kontrolliert wird, verschwindet der Beitrag der Benutzerfreundlichkeit zur Erklärung der Nutzung nahezu vollständig, in der ersten Untersuchung um 98 Prozent und in der zweiten um 91 Prozent. Davis liest das als Kette von der Benutzerfreundlichkeit über die Nützlichkeit zur Nutzung: Die Benutzerfreundlichkeit ist demnach eher eine Vorbedingung der wahrgenommenen Nützlichkeit als eine zweite, unabhängig wirkende Größe. Er selbst nennt diese Deutung vorläufig und fordert weitere Untersuchungen, was ihr Gewicht nicht aufhebt, aber ihre Reichweite begrenzt. Ebenso deutlich formuliert Davis die Gegenrichtung: Nutzende nehmen eine gewisse Bedienungsschwierigkeit in Kauf, wenn ein System dringend benötigte Funktionen bietet, während kein Maß an Benutzerfreundlichkeit ein System aufwiegt, das nichts Nützliches leistet.

\abbildung{abb-3-1-tam.png}{0.9}{abb:tam-pfade}%
{Die von Davis (1989) berichteten Zusammenhänge und die von Shroff et al. (2011) geprüften Pfade}%
{Die von Davis (1989) berichteten Zusammenhänge und die von Shroff et al. (2011) geprüften Pfade (eigene Darstellung nach Davis (1989) und Shroff et al. (2011))}

Abbildung 3.1 stellt beide Untersuchungen nebeneinander. Bei Davis führt der Weg von der wahrgenommenen Benutzerfreundlichkeit zur Nutzung über die Nützlichkeit. Bei Shroff et al. sind von vier geprüften Pfaden allein die beiden abgesichert, die von der wahrgenommenen Benutzerfreundlichkeit ausgehen; die Anschlüsse an die Nutzungsabsicht bleiben offen. Der gemeinsame Nenner ist damit schmaler, als er auf den ersten Blick wirkt, liegt aber genau an der Stelle, auf die es hier ankommt.

Shroff et al. (2011) übertragen das Modell auf den einschlägigen Gegenstand. Untersucht wurde die Absicht Studierender, ein E-Portfolio-System zu nutzen, das sie ein Semester lang eingesetzt hatten; von 169 ausgegebenen Fragebogen waren 72 auswertbar. Der Pfad von der wahrgenommenen Benutzerfreundlichkeit zur wahrgenommenen Nützlichkeit ist mit 0,71 der stärkste abgesicherte Zusammenhang des Modells. Für diesen einen Pfad zeigt eine zweite, unabhängige Erhebung damit in dieselbe Richtung wie Davis, und zwar an einem System der hier betrachteten Art. Eine Mediation prüfen Shroff et al. nicht.

Die Grenzen dieser Bestätigung gehören dazu. Die Stichprobe ist mit 72 Personen klein, sie stammt aus einer einzigen Einrichtung, und die Daten sind aus dem Frühjahr 2010. Der Einzelbefund zum Zusammenhang von Benutzerfreundlichkeit und Nützlichkeit ist dennoch belastbar, weil er nicht allein steht; die Wirkungsrichtung setzen beide Modelle voraus, prüfen kann sie keines von beiden.

Für die vorliegende Arbeit folgt daraus eine Festlegung, die weit vor der Evaluation greift. Ist die wahrgenommene Benutzerfreundlichkeit der Hebel auf die wahrgenommene Nützlichkeit, dann entscheidet sie darüber, ob ein freiwillig gefülltes Portal überhaupt gefüllt wird, und sie ist eine Frage des Entwurfs und nicht der Nacharbeit. Was das Modell nicht liefert, benennt Davis als offene Forschungsfrage: Wie sich Nützlichkeit und Benutzerfreundlichkeit beeinflussen lassen, bleibt bei ihm ungeklärt; als Ansatzpunkte nennt er Funktions- und Oberflächeneigenschaften, Entwicklungsvorgehen, Schulung und die Beteiligung der Nutzenden am Entwurf. Das Akzeptanzmodell begründet also, warum Gebrauchstauglichkeit zählt, und liefert ein geprüftes Messinstrument; welche Merkmale sie herstellen, beantwortet der folgende Abschnitt.

\section{Usability und Gestaltungsprinzipien}
\label{sec:usability}

Gebrauchstauglichkeit ist keine Eigenschaft, die einem System für sich zukommt. Brooke (1996) hält fest, dass sie sich nur mit Bezug auf einen bestimmten Nutzungskontext bestimmen lässt, und sieht diese Sicht im damaligen Normentwurf ISO 9241-11 gespiegelt; von dort übernimmt er die drei Klassen, die Maße der Gebrauchstauglichkeit abdecken sollen: Effektivität, Effizienz und Zufriedenstellung. Ein absolutes Maß kann es deshalb nicht geben, und den Vergleich von Systemen mit verschiedenem Zweck nennt Brooke einen zu vermeidenden Fall von Äpfeln und Birnen. Wenn überhaupt etwas systemübergreifend vergleichbar ist, dann nach seiner Einschätzung das Urteil derer, die ein System benutzen. Die Kriterien dieses Abschnitts sind damit an den Kontext eines institutionellen Portfolio-Portals gebunden.

Das Verfahren, mit dem sich eine Oberfläche ohne Testpersonen beurteilen lässt, geht auf Nielsen und Molich (1990) zurück. Sie ordnen die Möglichkeiten der Bewertung in vier Klassen: formal über eine Analysetechnik, automatisch über eine Prozedur, empirisch über Versuche mit Nutzenden und heuristisch, indem jemand die Oberfläche ansieht und nach eigenem Urteil bewertet. Diese vierte, informelle Klasse machen die Autoren arbeitsfähig, indem sie mehrere Begutachtende unabhängig voneinander an eine kleine Prinzipienliste binden; gängige Richtliniensammlungen liegen dagegen in der Größenordnung von tausend Regeln. Die Untersuchung von 1990 arbeitet mit neun Prinzipien. Die heute gebräuchliche Liste umfasst zehn: Aus einer Faktorenanalyse von 249 Usability-Problemen leitete Nielsen (1994a) neun Heuristiken ab, ergänzt um „Hilfe und Dokumentation“; in dieser Form wird die Liste seither unverändert geführt (Nielsen, 2024). Wer die Zehnerliste der Untersuchung von 1990 zuschreibt, zitiert deshalb falsch.

Vier Versuchsreihen mit 148 Begutachtenden, die zusammen 182 Einzelbegutachtungen erbrachten, zeigen, wie viel eine einzelne Begutachtung leistet: Je nach Oberfläche fand eine Person im Mittel zwischen 20 und 51 Prozent der bekannten Probleme. Werden mehrere unabhängige Befunde zusammengeführt, steigt der Anteil deutlich; von fünf Personen erwarten die Autoren rund zwei Drittel. Daraus stammt die Empfehlung, mit drei bis fünf Begutachtenden zu arbeiten. Ihre zweite Hälfte ist für die Anlage der Erhebung in Kapitel 6 folgenreich: Was darüber hinaus an Mitteln verfügbar ist, soll nicht in weitere Begutachtende fließen, sondern in andere Verfahren. Die Grenzen benennen die Autoren selbst: Ihr Verfahren zeigt Probleme mitunter ohne Lösungsvorschlag, es ist durch die Denkweise der Begutachtenden verzerrt, und geprüft wurde es allein an kleineren Oberflächen. Ergänzt werden soll es durch nutzerzentrierte Verfahren wie das laute Denken; dafür ist Nielsen (1994b) die methodische Referenz, auf der Kapitel 6 aufsetzt.

Die zehn Prinzipien heißen Heuristiken, weil sie breite Faustregeln sind und keine spezifischen Einzelrichtlinien (Nielsen, 2024); sie liefern Kriterien für eine Begutachtung und keine Abnahmekriterien. Tabelle 3.1 führt sie mit ihrer Kennung auf und nennt je Heuristik, was aus ihr für ein institutionelles Portfolio-Portal folgt. Die Spalte „Abgeleitetes Kriterium“ ist eine Ableitung der vorliegenden Arbeit und setzt den Zuschnitt eines solchen Portals voraus; sie ist der Maßstab für Abschnitt 7.2, während die Analyse in Abschnitt 3.4 nach den Heuristiken selbst bewertet.

\begingroup
\small
\setlength{\tabcolsep}{4pt}
\begin{longtable}{@{}L{1.7cm}L{4.4cm}L{8.3cm}@{}}
	\caption[Die zehn Usability-Heuristiken und das abgeleitete Kriterium]{Die zehn Usability-Heuristiken und das je Heuristik abgeleitete Kriterium für ein institutionelles Portfolio-Portal (eigene Darstellung nach Nielsen (2024))}
	\label{tab:heuristiken}\\
	\toprule
	\textbf{Kennung} & \textbf{Heuristik} & \textbf{Abgeleitetes Kriterium} \\
	\midrule
	\endfirsthead
	\toprule
	\textbf{Kennung} & \textbf{Heuristik} & \textbf{Abgeleitetes Kriterium} \\
	\midrule
	\endhead
	\bottomrule
	\endfoot
	H1 & Sichtbarkeit des Systemzustands & Der Zustand eines Vorgangs und die wirkende Einschränkung sind ablesbar \\
	H2 & Übereinstimmung von System und realer Welt & Die Benennungen folgen dem Vokabular der Hochschule, nicht dem des Systems \\
	H3 & Kontrolle und Freiheit der Nutzenden & Eine Einschränkung lässt sich einzeln und vollständig zurücknehmen \\
	H4 & Konsistenz und Standards & Gleichartige Inhalte folgen demselben Aufbau, wiederkehrende Elemente den Gewohnheiten des Web \\
	H5 & Fehlervermeidung & Folgenreiche Schritte werden verhindert, bevor sie erklärt werden müssen \\
	H6 & Wiedererkennen statt Erinnern & Die Ordnungsmerkmale des Bestandes werden sichtbar angeboten, statt Kenntnis vorauszusetzen \\
	H7 & Flexibilität und Effizienz der Nutzung & Wer den Bestand kennt, erreicht eine Arbeit ohne den Weg über alle Ordnungsebenen \\
	H8 & Ästhetisches und minimalistisches Design & Die Gestaltung ordnet sich den Arbeiten unter \\
	H9 & Unterstützung beim Erkennen, Verstehen und Beheben von Fehlern & Eine zurückgewiesene Eingabe nennt Grund und nächsten Schritt \\
	H10 & Hilfe und Dokumentation & Wer einen Vorgang erstmals durchläuft, findet die Erklärung im System \\
\end{longtable}
\endgroup

Die Auffindbarkeit, nach der die zweite Teilfrage fragt, führen die Heuristiken nicht als eigenes Thema. Auch die System Usability Scale (SUS), die weiter unten beschrieben wird, erfasst sie nicht: Lewis und Sauro (2018) halten fest, dass Eigenschaften wie die Auffindbarkeit in ihren zehn Aussagen nicht vorkommen. Die vorliegende Arbeit leitet sie deshalb aus drei Kriterien ab. Aus H6 folgt, dass die Systematik eines Bestandes als auswählbare Merkmale vorliegen muss, statt dass Suchende sie kennen müssen. Aus H3 folgt, dass eine Einschränkung einzeln und vollständig rücknehmbar sein muss. Aus H1 folgt, dass erkennbar bleiben muss, welche Einschränkungen wirken, weil sich ein leeres Ergebnis sonst nicht von einem leeren Bestand unterscheiden lässt. Dass die drei zusammen die Auffindbarkeit ergeben, ist eine Setzung dieser Arbeit.

Neben die auslegungsbedürftigen Heuristiken tritt eine Quelle anderer Art. Die Web Content Accessibility Guidelines 2.2 (WCAG 2.2) des World Wide Web Consortium (W3C, 2023) sind kein Forschungsbeitrag, sondern ein Webstandard mit überprüfbaren Mindestanforderungen, von denen ein Teil rechnerisch bestimmt ist. Für ein Portal mit öffentlichem Bestand und redaktionellem Ablauf sind zwölf Erfolgskriterien einschlägig; diese Auswahl trifft die vorliegende Arbeit, der Standard kennt keine solche Rangfolge. Der Textkontrast muss mindestens 4,5 zu 1 betragen, bei großer Schrift ab 18 Punkt oder 14 Punkt fett genügen 3 zu 1 (1.4.3); Bedienelemente und bedeutungstragende Grafik brauchen 3 zu 1 (1.4.11). Information darf nicht allein über Farbe vermittelt werden (1.4.1), was die Anzeige eines Bearbeitungszustands unmittelbar betrifft. Text muss sich auf 200 Prozent vergrößern lassen (1.4.4), bei 320 CSS-Pixeln Breite umbrechen (1.4.10) und größere Abstände vertragen (1.4.12). Der Tastaturfokus muss sichtbar sein (2.4.7) und darf nicht vollständig verdeckt werden (2.4.11), und Ziele für Zeigereingaben messen vorbehaltlich benannter Ausnahmen mindestens 24 mal 24 CSS-Pixel (2.5.8). Hinzu kommen drei Kriterien zu Bedienung und Auszeichnung: Alle Funktionen müssen per Tastatur erreichbar sein (2.1.1), Struktur und Beziehungen müssen maschinell erkennbar ausgezeichnet sein (1.3.1), und Bedienelemente müssen Name, Rolle und Zustand für Hilfsmittel bereitstellen (4.1.2). Die Kriterien 1.3.1, 1.4.1, 2.1.1 und 4.1.2 liegen auf der Stufe A, die übrigen acht auf der Stufe AA, dem Zielniveau dieser Arbeit. Barrierefreiheit ist damit kein Sonderthema neben der Benutzerfreundlichkeit, sondern ihr messbarer Teil.

Für die spätere Messung tritt ein drittes Instrument hinzu. Der Fragebogen von Brooke (1996), die SUS, besteht aus zehn abwechselnd zustimmend und ablehnend formulierten Aussagen; die Abwechslung zwingt die Befragten, jede einzeln zu bedenken. Aus den Antworten entsteht ein Summenwert zwischen 0 und 100. Dass dieser Wert trotz seiner Spanne kein Prozentwert ist, hält Brooke nicht ausdrücklich fest; die vorliegende Arbeit behandelt ihn als Punktwert in Schritten von 2,5. Brooke ordnet den Bogen als grobes Verfahren ein und nennt keinen Vergleichswert; ein einzelnes Ergebnis ist damit ohne Bezugsgröße nicht deutbar. Diese Bezugsgröße liefert eine Notenskala aus 241 industriellen Untersuchungen (Sauro \& Lewis, 2016, zitiert nach Lewis \& Sauro, 2018): Ein Wert von 68 liegt darauf in der Mitte der Note C und, da diese Mitte dem 50. Perzentil entspricht, zugleich im Median der Studienmittelwerte. Daraus stammt der in den Abschnitten 1.2 und 2.4 vorab gesetzte Zielwert. Brooke warnt davor, einzelne Aussagen für sich zu deuten; aus demselben Grund ordnet diese Arbeit sie den Konstrukten des Akzeptanzmodells nicht zu.

Ein Befund aus der Akzeptanzliteratur gehört an diese Stelle, weil er ein Gestaltungskriterium berührt. Davis (1989) berichtet, dass die Aussage zur Flexibilität für eine der untersuchten Anwendungen mit zwei der übrigen Aussagen zur Benutzerfreundlichkeit nicht signifikant zusammenhing, und folgert entgegen der landläufigen Annahme, dass Flexibilität nicht immer mit Benutzerfreundlichkeit einhergeht. Für ein Werkzeug, das gestalterische Freiheit einräumen soll, ist das ein unbequemer Befund aus einer Hauptquelle dieser Arbeit.

\section{Verwandte Arbeiten und Referenzplattformen}
\label{sec:verwandte-arbeiten}

Für den Entwurf eines hochschulischen Portfolio-Portals liegt wenig vor, an dem er sich ausrichten ließe. Schon für die schmalere Frage nach einem Inhaltsmodell des studentischen E-Portfolios sieht Mahasneh (2020) keine Vorarbeit; auf Portale bezogen ist diese Lücke eine Folgerung der vorliegenden Arbeit. Was die Literatur dagegen ordnet, sind die Wege zu einem solchen System. Lorenzo und Ittelson (2005) unterscheiden vier: Eigenentwicklung, quelloffene Systeme, kommerzieller Einkauf und der Rückgriff auf gängige Werkzeuge. Für die Eigenentwicklung führen sie drei Vorteile an, darunter, dass die Einrichtung genau das erhält, was sie will, und vier Nachteile, die alle dieses Vorhaben treffen: hohe Entwicklungskosten, hoher Zeitaufwand, hohe Anforderungen an Fachkompetenz und das Risiko, das nötige Personal nicht zu binden. Abschnitt 8.2 nimmt das wieder auf.

An derselben Stelle führen Lorenzo und Ittelson einen Vor- und einen Nachteil auf, aus denen sich ein Zielkonflikt lesen lässt. Gängige Werkzeuge erlauben freiere Portfolios, weil sie ihre Nutzenden nicht in feste Felder zwingen, verlangen dafür aber Kenntnisse im Erstellen von Webseiten; als Gegenbild nennen sie einen vorlagengestützten Weg, der diese Kenntnisse entbehrlich macht. Dass beides dieselbe Unterscheidung ist, ist eine Lesart dieser Arbeit. Shroff et al. (2011) halten die eine Seite des Konflikts fest: Das Gestaltungsmodell bestehender E-Portfolio-Systeme könne begrenzt und starr darin sein, wie Fluss und Erscheinungsbild der Inhalte verwaltet werden. Der Satz steht in ihrem Forschungsstand und ist kein Ergebnis ihrer Erhebung. Gestalterische Freiheit ohne Kenntnisse im Erstellen von Webseiten ist damit keine Vorliebe, sondern die Antwort auf ein benanntes Problem.

Eine Plattform, die dem Vorhaben im Anwendungsfall nahekommt, ist das Portfolio-Portal der Hochschule für Gestaltung Schwäbisch Gmünd. Es zeigt die Arbeiten aus vier Bachelor-Studiengängen ohne Anmeldung und wendet sich damit wie das Vorhaben an Nutzende ohne Konto. Herangezogen wurde es auf Empfehlung des Erstprüfers und wegen dieser Vergleichbarkeit; ein systematisches Auswahlverfahren liegt der Wahl nicht zugrunde, und eine zweite Plattform wurde nicht untersucht. Die Befunde stützen deshalb keine Aussage über hochschulische Portfolio-Portale allgemein.

Die Analyse ist eigene Vorarbeit und keine Literatur. Sie wurde im Mai 2026 von den beiden Verfassern dieser Arbeit durchgeführt, nach dem Verfahren von Nielsen und Molich (1990) und mit der Zehnerliste aus Tabelle 3.1 als Raster. Begutachtet wurden fünf Klassen von Seiten, von der Startseite über die Studiengangsübersichten bis zu den Schlagwort- und Personenansichten. Bewertet wurde vierstufig; neben „erfüllt“, „teilweise erfüllt“ und „nicht erfüllt“ steht die Stufe „nicht beurteilbar“ für alles hinter der Anmeldung, denn ein nicht einsehbarer Bereich ist kein Mangel. Tabelle 3.2 fasst den Befundstand zusammen.

\begingroup
\small
\setlength{\tabcolsep}{4pt}
\begin{longtable}{@{}L{1.7cm}L{3.1cm}L{9.6cm}@{}}
	\caption[Befundstand der heuristischen Analyse der Referenzplattform]{Befundstand der heuristischen Analyse der Referenzplattform (eigene Darstellung, heuristische Analyse Mai 2026)}
	\label{tab:hfg-befunde}\\
	\toprule
	\textbf{Kennung} & \textbf{Bewertung} & \textbf{Kernbefund} \\
	\midrule
	\endfirsthead
	\toprule
	\textbf{Kennung} & \textbf{Bewertung} & \textbf{Kernbefund} \\
	\midrule
	\endhead
	\bottomrule
	\endfoot
	H1 & teilweise erfüllt & Keine Brotkrumennavigation, aktive Auswahl nur im Seitenkopf und im Bedienelement \\
	H2 & erfüllt & Sprache und Modulbezeichnungen folgen dem Curriculum \\
	H3 & teilweise erfüllt & Kein sichtbares Zurücksetzen der Auswahl \\
	H4 & erfüllt & Alle Projektseiten folgen demselben Aufbau, Adressen einheitlich \\
	H5 & nicht beurteilbar & Sichtbare Schreibvorgänge unauffällig, folgenreiche Pfade hinter der Anmeldung \\
	H6 & teilweise erfüllt & Keine Gesamtübersicht der Schlagworte \\
	H7 & nicht erfüllt & Kein Suchfeld, Filter nicht kombinierbar \\
	H8 & erfüllt & Zurückhaltende Gestaltung, die Projekte tragen die Seite \\
	H9 & nicht beurteilbar & Einreichungsweg nicht einsehbar, Fehlerseiten nicht systematisch geprüft \\
	H10 & teilweise erfüllt & Kurzer Einführungstext, keine Hilfe- oder Fragenseite \\
\end{longtable}
\endgroup

Drei Heuristiken sind erfüllt, vier teilweise, zwei blieben unbeurteilt, und eine ist nicht erfüllt. Bei H5 tragen zwei Beobachtungen die Einstufung: Die sichtbaren Schreibvorgänge blieben unauffällig, und die folgenreichen Pfade liegen hinter der Anmeldung. Bei H9 kommt hinzu, dass Fehlerseiten nicht systematisch geprüft wurden; das ist eine Lücke der eigenen Erhebung und keine Zugangsbeschränkung. Der einzige Negativbefund liegt bei H7 und trifft die Hauptzielgruppe unmittelbar.

\abbildung{abb-3-2-hfg-filter.png}{0.95}{abb:hfg-uebersicht}%
{Studiengangsübersicht der Referenzplattform}%
{Studiengangsübersicht der Referenzplattform: eine Auswahlliste als vollständiger Filterapparat (Bildschirmaufnahme portfolio.hfg-gmuend.de, abgerufen am 05.09.2026)}

Abbildung 3.2 zeigt eine der vier Studiengangsübersichten; die Aufnahme entstand vier Monate nach der Begehung und zeigt einen unveränderten Stand. Unter dem Seitentitel steht eine einzige Auswahlliste, und sie ist der gesamte Filterapparat; ein Suchfeld gibt es nicht, eine Brotkrumennavigation ebenso wenig, und welche Einschränkung wirkt, sagen allein der Seitenkopf und das Bedienelement selbst. Auf den Schlagwort- und Personenansichten fehlt auch dieses. Die Projekte tragen dagegen Modul und Semester in der Kachel, was H2 und H4 stützt: Die Benennungen stammen aus dem Curriculum, und alle Kacheln sind gleich aufgebaut.

Über die einzelnen Heuristiken hinaus ergeben sich vier Beobachtungen zur Informationsarchitektur; eine Nachprüfung am 05.09.2026 hat die erste Begehung dabei in drei Punkten berichtigt. Erschlossen ist der Bestand über Studiengang, Schlagwort und Person. Semester und Modul erscheinen zwar an jedem Projekt, sind aber keine auswählbaren Achsen: Das Semesterkürzel in der Adresse schränkt die Trefferliste nicht ein, und eine Modulansicht gibt es nicht. Die mehrachsige Erschließung ist damit nur zum Teil ausgebildet, und die vorhandenen Achsen lassen sich ohnehin nicht verbinden. Projekt- und Studiengangsansichten tragen sprechende, dauerhafte Adressen; Schlagwort- und Personenansichten sind über undurchsichtige Kurzkennungen adressiert. Die Projektseiten sind erzählend angelegt und nicht als Datenblatt. Studierende erscheinen bei neueren Projekten als verlinkte Sammelseiten ihrer Arbeiten, allerdings ohne Biografie oder Bild, und Betreuende stehen als reiner Text. Die Person ist hier eine Erschließungsachse und noch kein eigener Gegenstand.

Die Reichweite dieser Befunde ist begrenzt. Die Analyse wurde von zwei Personen durchgeführt und bleibt damit unter den drei bis fünf, die Nielsen und Molich (1990) empfehlen; nach deren eigenen Zahlen ist dann nur ein Teil der vorhandenen Probleme zu erwarten. Beurteilt wurde ausschließlich die Sicht nicht angemeldeter Besuchender, die Begehung erfolgte stichprobenartig, und weder Ladezeiten noch die in Abschnitt 3.3 genannten Erfolgskriterien wurden geprüft. Hinzu kommen die Grenzen, die Nielsen und Molich dem Verfahren zuschreiben. Die Befunde sind deshalb Hinweise und keine Mängelliste. Da die Außensicht zugleich die Sicht der Hauptzielgruppe ist, mindert die Beschränkung auf die Sicht ohne Anmeldung die Aussagekraft für die Auffindbarkeit im öffentlichen Bereich nicht; für den Einreichungsweg fehlt dagegen jede Vorlage.

\section{Zusammenfassung}
\label{sec:zusammenfassung-grundlagen}

\platzhalter{Abschnitt 3.5 steht auf der Confluence-Seite 34373645 noch nicht. Er beantwortet die erste Teilfrage aus der Forschungsliteratur und fasst zusammen, was die drei Literaturstränge jeweils tragen.}
